% =============================================================================
% Labs/CS301/06-linked-lists-singly-coding/06-linked-lists-singly-coding-assignment.tex
% Coding Lab 6: The Roster That Rewires Itself (singly linked lists)
%
% Student-facing problem statement. Instructor-only files live in this same
% directory but are NEVER published to students or synced to docs/:
%   reference_solution.c  (NEVER published)
%   test_hidden.c         (NEVER published)
%   autograder.sh         (NEVER published)
% Student-facing: slist.h, starter_slist.c, test_visible.c (+ this PDF).
%
% Fresh instances throughout: no worked example re-asks a deck/notes,
% written-assignment, or earlier-lab instance verbatim (deck/notes: build
% 10-20-30; front-insert 5 onto head->10->20; position insert 25 at pos 2;
% delete 10 head + delete 30 tail; search/display on 10-20-30 and 5-15-25;
% reverse 1-2-3 and 10-20-30; before-you-go head->5->15 insert 10, delete 15,
% reverse 1-2-3; written Q3: 35-40-50-55-60 chain, Q4: 7-14-21-28, Q5:
% 11-22-29-33-44 with x=29/5/50, Q6: 6-12-18-24). Visible tests use the
% roster chain 17-42-88 with insert 63 and absent-lookup 77 — all new.
%
% Compile with: cd Labs/CS301/06-linked-lists-singly-coding &&
%   TEXINPUTS=../../../Preambles;$TEXINPUTS xelatex -interaction=nonstopmode 06-linked-lists-singly-coding-assignment.tex
%   (Windows/MiKTeX uses ; not : as the path separator)
% =============================================================================

\documentclass[11pt]{article}
\usepackage[margin=1in]{geometry}
% =============================================================================
% Preambles/header.tex — shared LaTeX/Beamer preamble for this project
%
% Use from a Beamer lecture by prepending:
%     \input{header}
% and compiling with TEXINPUTS=../Preambles:$TEXINPUTS (the /compile-latex
% skill does this automatically).
%
% PALETTE CONTRACT. Color names defined here MUST match the names in
% Quarto/theme-template.scss so Beamer and Quarto renderings use the same
% palette. The scripts/check-palette-sync.sh script enforces this.
%
% To customize: edit the HEX values below AND the matching $variable in
% Quarto/theme-template.scss, then run scripts/check-palette-sync.sh.
% =============================================================================

% -----------------------------------------------------------------------------
% Engine detection + encoding
%
% Our /compile-latex skill uses XeLaTeX, where inputenc is unnecessary and
% can produce encoding warnings. Only load inputenc under pdfTeX.
% -----------------------------------------------------------------------------
\usepackage{iftex}
\ifPDFTeX
  \usepackage[utf8]{inputenc}
\fi

\usepackage{amsmath, amssymb, amsthm}
\usepackage{booktabs}
\usepackage{graphicx}

% -----------------------------------------------------------------------------
% PALETTE — mirrors Quarto/theme-template.scss variable names.
% Load xcolor and define colors BEFORE anything that references them
% (notably hyperref, hypersetup, and the Beamer color assignments).
% -----------------------------------------------------------------------------
\usepackage{xcolor}

\definecolor{primary-blue}{HTML}{012169}      % headings, accents
\definecolor{primary-gold}{HTML}{B9975B}      % emphasis, borders
\definecolor{highlight-yellow}{HTML}{F2A900}  % markers, alerts
\definecolor{light-bg}{HTML}{E8EDF5}          % title / section backgrounds
\definecolor{jet}{HTML}{1A1A1A}               % body text

% Semantic colors (used in diagrams and annotations)
\definecolor{positive}{HTML}{15803D}          % good / observed / identified
\definecolor{negative}{HTML}{B91C1C}          % bad / problematic / bias
\definecolor{neutral}{HTML}{525252}           % reference / context

% Highlight accents (match .hi-* classes in the SCSS)
\definecolor{hi-slate}{HTML}{314F4F}
\definecolor{hi-green}{HTML}{2E8B57}
\definecolor{hi-red}{HTML}{C41E3A}

% Semantic aliases so skill templates and snippets can write
% \textcolor{accent}{...} without hard-coding "primary-blue".
\colorlet{accent}{primary-blue}
\colorlet{accent2}{primary-gold}

% -----------------------------------------------------------------------------
% Hyperref — loaded AFTER xcolor + palette so link colors resolve.
% -----------------------------------------------------------------------------
\usepackage{hyperref}
\hypersetup{colorlinks=true, linkcolor=primary-blue, urlcolor=primary-blue,
            citecolor=primary-blue, pdfborder={0 0 0}}

% -----------------------------------------------------------------------------
% Course code — set once per deck (after \input{header}, before \begin{document})
% via \coursecode{CS401}. Shown in the footer on every slide so a deck's course
% is identifiable without opening the title page. Empty by default.
% -----------------------------------------------------------------------------
\makeatletter
\newcommand{\coursecode}[1]{\gdef\@coursecodetext{#1}}
\gdef\@coursecodetext{}
\makeatother

% -----------------------------------------------------------------------------
% Beamer theme assignments (only apply under Beamer).
%
% We detect Beamer via \@ifclassloaded{beamer}, which is the documented,
% non-internal way. This must be wrapped in \makeatletter/\makeatother
% because \@ifclassloaded contains an @-catcode character.
% -----------------------------------------------------------------------------
\makeatletter
\@ifclassloaded{beamer}{%
  \usefonttheme{professionalfonts}
  \setbeamercolor{structure}{fg=primary-blue}
  \setbeamercolor{frametitle}{fg=primary-blue}
  \setbeamercolor{title}{fg=primary-blue}
  \setbeamercolor{subtitle}{fg=primary-gold}
  \setbeamercolor{author}{fg=primary-blue}
  \setbeamercolor{institute}{fg=neutral}
  \setbeamercolor{date}{fg=primary-gold}
  \setbeamercolor{itemize item}{fg=primary-blue}
  \setbeamercolor{itemize subitem}{fg=primary-gold}
  \setbeamercolor{alerted text}{fg=primary-gold}
  \setbeamercolor{block title}{fg=white, bg=primary-blue}
  \setbeamercolor{block body}{bg=light-bg}
  % Semantic block variants: block=definition/concept (blue), exampleblock=
  % worked example (green/positive), alertblock=key takeaway or caution (gold).
  % Keep to <=2 colored boxes per slide across ALL three variants (INV-7).
  \setbeamercolor{block title example}{fg=white, bg=positive}
  \setbeamercolor{block body example}{bg=light-bg}
  \setbeamercolor{block title alerted}{fg=white, bg=primary-gold}
  \setbeamercolor{block body alerted}{bg=light-bg}

  % Minimal footer: course code (left) + page X / N (right), no navigation clutter
  \setbeamertemplate{navigation symbols}{}
  \setbeamertemplate{footline}{%
    \color{neutral}\footnotesize\hspace{0.7em}\@coursecodetext\hfill
    \insertframenumber/\inserttotalframenumber\hspace{0.7em}\vspace{0.3em}}
}{}
\makeatother

% -----------------------------------------------------------------------------
% TikZ setup — libraries the snippet gallery uses
% -----------------------------------------------------------------------------
\usepackage{tikz}
\usetikzlibrary{arrows.meta, positioning, calc, decorations.pathreplacing,
                fit, shapes.geometric, backgrounds}

% Shared TikZ styles — snippets and hand-written diagrams can pick these up
% via `\begin{tikzpicture}[every node/.style={...}]` or by naming specific
% styles (e.g., `\node[dag-node] (X) at (0,0) {X};`).
\tikzset{
  % Node styles for causal DAGs and similar
  dag-node/.style = {
    draw=primary-blue, fill=light-bg, rounded corners=2pt,
    minimum width=1.2cm, minimum height=0.9cm, align=center, font=\small
  },
  decision-node/.style = {
    draw=primary-gold, fill=white, diamond, aspect=1.8,
    minimum width=1.8cm, minimum height=1.2cm, align=center, font=\small,
    inner sep=1pt
  },
  % Edge styles — observed vs counterfactual
  observed-edge/.style = {-{Latex[length=2.5mm]}, thick, draw=jet},
  counterfactual-edge/.style = {-{Latex[length=2.5mm]}, thick, draw=neutral, dashed},
  confound-edge/.style = {-{Latex[length=2.5mm]}, thick, draw=negative, dashed},
  % Data-point styles for plots
  observed-dot/.style = {fill=primary-blue, draw=primary-blue, circle, inner sep=1.2pt},
  counterfactual-dot/.style = {fill=white, draw=neutral, circle, inner sep=1.2pt}
}

% -----------------------------------------------------------------------------
% Convenience macros
% -----------------------------------------------------------------------------
\newcommand{\muted}[1]{\textcolor{neutral}{#1}}
\newcommand{\key}[1]{\textcolor{primary-gold}{\textbf{#1}}}
\newcommand{\good}[1]{\textcolor{positive}{#1}}
\newcommand{\bad}[1]{\textcolor{negative}{#1}}

% Standout frame macro: full-bleed dark background for section transitions.
% Beamer-only (uses \begin{frame}/\setbeamercolor/\usebeamerfont) -- do not
% call from an article-class file (e.g. Notes/<CODE>/*-notes.tex). \muted,
% \key, \good, \bad, and \coursecode above are plain xcolor/gdef and are
% safe to use from either class; verified when Notes/article-class support
% was added.
% Expand: \transitionslide{Your transition title}
\newcommand{\transitionslide}[1]{%
  \begingroup
  \setbeamercolor{background canvas}{bg=primary-blue}
  \begin{frame}[plain]
    \centering
    \vfill
    {\usebeamerfont{title}\color{white}\Large #1\par}
    \vfill
  \end{frame}
  \endgroup
}

% Section divider frame (Beamer-only), an alternative to \transitionslide with a
% darker two-line style: near-black (jet) background, a small white label line
% above a larger gold title, and a thin gold rule along the bottom edge.
% Expand: \sectiondivider{Section 2}{Pointers}
\newcommand{\sectiondivider}[2]{%
  \begingroup
  \setbeamercolor{background canvas}{bg=jet}
  \begin{frame}[plain]
    \vspace*{3.0cm}
    {\color{white}\ttfamily\large #1\par}
    \vspace{0.5cm}
    {\color{primary-gold}\LARGE #2\par}
    \begin{tikzpicture}[remember picture, overlay]
      \draw[primary-gold, line width=2.5pt]
        ($(current page.south west)+(0,0.1)$) -- ($(current page.south east)+(0,0.1)$);
    \end{tikzpicture}
  \end{frame}
  \endgroup
}

% -----------------------------------------------------------------------------
% Channel watermark — "Concepts That Click" (YouTube).
%
% Article-class files (CompetitiveExam Q/A sets, Notes, Assignments) get a
% light diagonal draft watermark on every page plus a centered footer naming
% the channel. Beamer decks get a subtle TikZ background watermark instead
% (the footline already carries the course code). Centralized here so every
% MCQ PDF is branded without per-file edits.
% -----------------------------------------------------------------------------
\makeatletter
\@ifclassloaded{article}{%
  \usepackage{draftwatermark}
  \SetWatermarkText{Concepts That Click}
  \SetWatermarkScale{1}
  \SetWatermarkFontSize{2cm}
  \SetWatermarkLightness{0.86}
  \SetWatermarkAngle{45}
  \usepackage{fancyhdr}
  \pagestyle{fancy}
  \fancyhf{}
  \fancyfoot[C]{\footnotesize\textcolor{neutral}{Concepts That Click~|~YouTube: \href{https://www.youtube.com/@concepts.thatclick}{@concepts.thatclick}}}
  \renewcommand{\headrulewidth}{0pt}
  \renewcommand{\footrulewidth}{0pt}
  \fancypagestyle{plain}{%
    \fancyhf{}%
    \fancyfoot[C]{\footnotesize\textcolor{neutral}{Concepts That Click~|~YouTube: \href{https://www.youtube.com/@concepts.thatclick}{@concepts.thatclick}}}%
    \renewcommand{\headrulewidth}{0pt}%
    \renewcommand{\footrulewidth}{0pt}%
  }
}{}
\@ifclassloaded{beamer}{%
  \setbeamertemplate{background}{%
    \begin{tikzpicture}[remember picture, overlay]
      \node[rotate=45, opacity=0.055, align=center]
        at (current page.center)
        {\Huge\textcolor{neutral}{Concepts That Click}};
    \end{tikzpicture}%
  }
}{}
\makeatother 
\coursecode{CS301}
\renewcommand{\thesection}{6.\arabic{section}}

\title{Coding Lab 6: The Roster That Rewires Itself}
\author{Auqib Hamid Lone\\[0.3em]
  \emph{Department of Computer Science \& Engineering}, \emph{GCET Ganderbal}}
\date{\today}

\begin{document}
\maketitle

\noindent\textbf{Due:} two weeks from issue. There is no automatic late penalty --- if you need more time, ask before the deadline and an extension will be granted (see the course policy in the syllabus).

\section{Why this problem}

Weeks 3--5 built an expression processor that can tokenize, check brackets, convert, evaluate, and buffer jobs in arrival order. One piece is still missing: evaluating \texttt{(a+b)*c} needs the \emph{values} of \texttt{a}, \texttt{b}, \texttt{c}, and identifiers arrive at unpredictable times in unpredictable numbers. A fixed array demands we guess the capacity up front and pay the slot tax on every front insertion. This lab is that missing piece: a roster of IDs held as a singly linked chain, where each entry carries the address of its successor so the lineup changes by rewriting one address instead of moving $n$ elements. You implement the five operations from lecture --- insertion (front, end, position), deletion, search, display, and in-place reverse --- with the same pointer discipline the deck drilled: save the only path \emph{before} you overwrite it, and free exactly once what you malloc'd.

\section{Your task}

Implement the singly linked list ADT in C on heap nodes with explicit contracts on every fallible operation. The contract is given to you in \texttt{slist.h} --- \emph{do not modify it}. Copy \texttt{starter\_slist.c} to \texttt{slist.c} and fill in all nine functions:

\begin{verbatim}
SNode  *slist_insert_front(SNode **head, int x);
SNode  *slist_insert_end(SNode **head, int x);
int     slist_insert_at(SNode **head, size_t pos, int x); /* 0 ok / -1 fail */
int     slist_delete_value(SNode **head, int x);          /* 0 gone / -1 absent */
SNode  *slist_search(SNode *head, int x);                 /* NULL when absent */
size_t  slist_length(const SNode *head);
int     slist_display(const SNode *head, int *out, int out_cap);
void    slist_reverse(SNode **head);
void    slist_free(SNode **head);
\end{verbatim}

\textbf{Insertion is the deck's walk-then-splice.} Front-insert is the two-write sequence (\texttt{newNode->next = *head} first, \emph{then} \texttt{*head = newNode}) --- swapped, the old chain is orphaned with no copy anywhere. End-insert walks to the last node (an empty list sets \texttt{*head} directly); position insert walks \texttt{p} to the predecessor then splices in the same order, refusing \texttt{pos} past the end. Every \texttt{malloc} is checked for \texttt{NULL} before dereferencing.

\textbf{Deletion is bypass, reclaim, never touch.} Delete the \emph{first} node holding \texttt{x} (head first, else the predecessor walk): rewrite one link, \texttt{free} exactly once, never touch the node afterwards. Absent \texttt{x} changes nothing and reports $-1$; only the first duplicate goes.

\textbf{Reverse is the three-pointer loop.} \texttt{prev = NULL}, \texttt{curr = *head}; save \texttt{next}, flip, advance both walkers; \texttt{*head = prev}. Display never disturbs the list (the traversal pointer moves while \texttt{head} stays put), and \texttt{free} releases every node and NULLs \texttt{*head} so a second free is safe.

\section{Constraints}

\begin{itemize}
  \item Representation must be the heap-chained \texttt{SNode} nodes of \texttt{slist.h} --- no arrays, no global/static mutable state.
  \item Every allocation is checked (\texttt{malloc} for \texttt{NULL}); every node is freed exactly once; no use-after-free, no writes past allocated bounds. Error paths leave the list unchanged.
  \item Correctness is graded against the contract; your code must compile with \texttt{clang -std=c11 -Wall -Wextra -Werror} without errors or warnings.
  \item Complexity: front-insert $O(1)$ worst case; end/position insert, delete, and search $O(n)$ worst case in the current length (search $O(1)$ best case when \texttt{head} holds \texttt{x}); display $\Theta(n)$; reverse $O(n)$ worst case in $O(1)$ auxiliary space; length and free $O(n)$ worst case.
\end{itemize}

\section{Worked examples (these are the visible tests)}

\textbf{Example 1 --- build the roster.} Front-insert $42$ (chain: $42$), front-insert $17$ (chain: $17 \to 42$), end-insert $88$ (chain: $17 \to 42 \to 88$). Length is $3$; display writes $17, 42, 88$. Freeing sets \texttt{head} to \texttt{NULL}.

\textbf{Example 2 --- splice, delete, search.} From $17 \to 42 \to 88$: insert $63$ at position $1$ giving $17 \to 63 \to 42 \to 88$; delete the first $42$ giving $17 \to 63 \to 88$ with length $3$; \texttt{search(63)} finds its node while \texttt{search(77)} reports absent (\texttt{NULL}).

\textbf{Example 3 --- reverse and edges.} Reversing $17 \to 63 \to 88$ gives $88 \to 63 \to 17$; deleting the head $88$ leaves length $2$; reversing again gives $17 \to 63$. Reversing the freed (empty) list is a safe no-op with length $0$.

\section{What to submit}

\begin{enumerate}
  \item \texttt{slist.c} --- your completed implementation (edit \texttt{starter\_slist.c}).
  \item \texttt{slist.h} --- unchanged, but include it so the submission compiles standalone.
\end{enumerate}

\section{Getting started}

\begin{enumerate}
  \item Copy \texttt{starter\_slist.c} to \texttt{slist.c} and fill in the nine function bodies.
  \item Sanity-check locally:
\begin{verbatim}
gcc -std=c11 -Wall -Wextra -o test_visible test_visible.c slist.c
./test_visible
\end{verbatim}
    All three worked examples should pass.
  \item Submit \texttt{slist.c} and \texttt{slist.h}. The autograder compiles your \texttt{slist.c} against a \emph{hidden} test suite and reports pass/fail. The hidden suite covers: \texttt{NULL} head-pointer safety for every op; front/end/position insert boundaries including \texttt{pos} $==$ length vs.\ past-the-end refusal with the list unchanged; head/middle/tail/absent/empty deletes with first-occurrence-only duplicates; search identity at head and tail; display validation with too-small, \texttt{NULL}, and negative-capacity buffers plus a canary check that nothing is written on refusal; exact-capacity display and head purity; \texttt{INT\_MIN}/\texttt{INT\_MAX}/\texttt{-1}/\texttt{0} round-trips; reverse of empty/single/two/four chains with double-reverse identity; free with head-NULLing and double-free safety; and a seeded 200-operation random campaign checked operation-by-operation against an independent array oracle.
\end{enumerate}

\end{document}
